%! AP Statistics — Unit 1, Lesson 1.3 — Homework ANSWER KEY
\documentclass[10pt]{article}
\usepackage{apstats-article}
\usepackage{apstats-key}

%=============================================================================
\begin{document}

\pageheader{Unit 1, Lesson 1.3}{Homework --- Answer Key}
\begin{tabularx}{\linewidth}{X X X}
Name:\;\ans{Key} & Date:\; & Period:\;
\end{tabularx}

\vspace{0.1in}

% ════════════════════════════════════════════════════════════════════════════
% PART A
% ════════════════════════════════════════════════════════════════════════════
\begin{blurbbox}[Part A \quad Build a Frequency Table from Raw Data]{navylight}
\small

Counts from the raw data: Classical~4, Hip-Hop~6, Pop~8, Rock~7 (total = 25).

\vspace{8pt}
\renewcommand{\arraystretch}{2.0}
\begin{tabularx}{\linewidth}{@{} >{\bfseries}p{0.22\linewidth}
                                    C{0.16\linewidth}
                                    C{0.2\linewidth}
                                    C{0.22\linewidth} @{}}
\rowcolor{sky}
\textbf{Music type} & \textbf{Tally} & \textbf{Frequency} & \textbf{Rel.\ Frequency} \\
\midrule
Classical  & \small{||||} & \ans{4} & \ans{0.16 (16\%)} \\
Hip-Hop    & \small{||||||} & \ans{6} & \ans{0.24 (24\%)} \\
Pop        & \small{||||||||} & \ans{8} & \ans{0.32 (32\%)} \\
Rock       & \small{|||||||} & \ans{7} & \ans{0.28 (28\%)} \\
\midrule
\rowcolor{sky}\textbf{Total} & & \ans{25} & \ans{1.00} \\
\end{tabularx}

\vspace{10pt}

\begin{enumerate}[leftmargin=1.5em, itemsep=8pt]
  \item \ans{Of the 25 students surveyed, pop was the most preferred music type,
        chosen by 32\% of students. Classical was the least popular, preferred by
        only 16\% of students. Hip-hop (24\%) and rock (28\%) were also popular
        choices.}

  \item \ans{Pop or rock: $0.32 + 0.28 = 0.60 = 60\%$ (or $(8+7)/25 = 15/25$).}
\end{enumerate}

\end{blurbbox}

\vspace{0.14in}

% ════════════════════════════════════════════════════════════════════════════
% PART B
% ════════════════════════════════════════════════════════════════════════════
\begin{blurbbox}[Part B \quad Compare Two Groups]{greenacc}
\small

\renewcommand{\arraystretch}{1.6}
\begin{tabularx}{\linewidth}{@{} >{\bfseries}p{0.18\linewidth}
                                    C{0.14\linewidth} C{0.2\linewidth}
                                    C{0.14\linewidth} C{0.2\linewidth} @{}}
\rowcolor{sky}
\textbf{Season} & \textbf{Class A} & \textbf{Rel.\ Freq.\ A}
                & \textbf{Class B} & \textbf{Rel.\ Freq.\ B} \\
\midrule
Spring & 6  & \ans{0.200 (20\%)} & 12 & \ans{0.300 (30\%)} \\
Summer & 14 & \ans{0.467 (46.7\%)} & 16 & \ans{0.400 (40.0\%)} \\
Fall   & 5  & \ans{0.167 (16.7\%)} & 10 & \ans{0.250 (25.0\%)} \\
Winter & 5  & \ans{0.167 (16.7\%)} & 2  & \ans{0.050 (5.0\%)} \\
\midrule
\rowcolor{sky}\textbf{Total} & 30 & \ans{1.001*} & 40 & \ans{1.000} \\
\end{tabularx}
\textcolor{slate}{\scriptsize *Rounding; accept 1.00 or 100\%.}

\vspace{10pt}

\begin{enumerate}[leftmargin=1.5em, itemsep=8pt]
  \item \ans{Summer is the most popular season in both classes. Class A: 46.7\%;
        Class B: 40.0\%.}

  \item \ans{No. Class B is a larger group (40 vs.\ 30). Using relative frequency:
        Class A has 46.7\% preferring summer vs.\ 40.0\% in Class B. Summer is
        actually more popular \textit{proportionally} in Class A.}

  \item \ans{Class A shows a stronger preference for winter: 16.7\% (5/30) vs.\
        5.0\% (2/40). Despite equal counts (both have 5 or fewer), the proportion
        in Class A is much higher.}

  \begin{teachernote}
  Class A winter = 5/30 $\approx$ 16.7\%; Class B winter = 2/40 = 5.0\%. This is the
  key point: same absolute count (or even Class A having more) can mean very
  different proportions.
  \end{teachernote}
\end{enumerate}

\end{blurbbox}

\vspace{0.14in}

% ════════════════════════════════════════════════════════════════════════════
% PART C
% ════════════════════════════════════════════════════════════════════════════
\begin{blurbbox}[Part C \quad AP-Style Free Response]{redacc}
\small

\begin{enumerate}[leftmargin=1.5em, itemsep=10pt]
  \item \ans{Of the 120 college athletes sampled, the dominant hand was right for
        the vast majority --- 80\% of athletes were right-handed. Left-handed
        athletes made up 15\% of the sample, and only 5\% were ambidextrous.}

  \item \ans{The table would be more informative if compared to the general
        population (where approximately 10\% of people are left-handed). Comparing
        relative frequencies across athlete and non-athlete groups could reveal
        whether left-handedness is more or less common among athletes.}

  \begin{teachernote}
  For (1): require $n = 120$, the variable, all three categories with relative
  frequencies, and the most/least common identified. Full AP credit requires
  contextual interpretation, not just reading the numbers.\\
  For (2): any reasonable comparison (athletes vs.\ non-athletes, different sports,
  different genders) earns credit. The key insight is that the table alone doesn't
  tell you whether 80\% right-handed is unusual or typical.
  \end{teachernote}
\end{enumerate}

\end{blurbbox}

\end{document}
